Theorem:
$$\sum_{k=1}^\infty\frac{H_{2k}-H_{k}}{k^2\binom{2k}k}=2 \zeta(3)-\frac{\pi  \sqrt{3}\, \Psi^{\left(1\right)}\! \left(\frac{1}{3}\right)}{24}+\frac{\pi  \sqrt{3}\, \Psi^{\left(1\right)}\! \left(\frac{5}{6}\right)}{72}\tag{1},$$
where $H_n$ denotes the [harmonic numbers][2] and $\Psi^{(n)}(x)$ is the $n$th [polygamma function][3].


Proof:
$\newcommand{\mfrac}[2]{{\textstyle\frac{#1}{#2}}}$

Let $S$ be the sum and let 
$$
F(z)=\sum_{k=1}^\infty \frac{z^k}{k^2\binom{2k}{k}} $$
for $0 \le z \le 1$. From the standard Taylor series
$$\arcsin^2 x = \frac{1}{2} \sum_{k=1}^\infty \frac{1}{k^2 \binom{2k}{k}} (2x)^{2k},$$
see for instance <a href="https://math.stackexchange.com/questions/1448822/taylor-expansion-for-arcsin2x">this MathStackexchange answer</a> (although ChatGPT was keen to share its own proof with me), we have $F(z)=2\arcsin^2\bigl(\frac{\sqrt z}{2}\bigr)$.
Using the integral formula for the Harmonic numbers we have
$$ H_{2k}-H_k
= \int_0^1 \frac{x^k-x^{2k}}{1-x}\,\mathrm{d}x $$
and so
$$ S =\int_0^1 \frac{F(x)-F(x^2)}{1-x}\,\mathrm{d}x. $$
Now
$F(x)-F(x^2)=\int_{x^2}^x F'(t)\,\mathrm{d}t $
so we have, swapping the order of integration using
$x^2 \le t \le x \iff t \le x \le \sqrt{t}$ to get the second line,
$$  \begin{aligned} S &= \int_0^1 \int_{x^2}^x F'(t)\, \mathrm{d}t\, \frac{\mathrm{d}x}{1-x} \\ &=
\int_0^1 F'(t)\bigl(\int_t^{\sqrt t}\frac{dx}{1-x}\bigr)\mathrm{d}t \\ &=
\int_0^1 F'(t)\log \bigl(\frac{1-t}{1-\sqrt t}\bigr)dt \\ 
&= \int_0^1 F'(t)\log(1+\sqrt t)\,\mathrm{d}t \\
&= \int_0^1 F'(u^2) \log (1+u)\, 2u\, \mathrm{d}u \\
&= \int_0^{\pi/6} F'(4\sin^2\theta) \log(1 + 2\sin \theta) 8\sin\theta \cos \theta \,\mathrm{d}\theta \\
&= \int_0^{\pi/6} \frac{\theta}{\sin 2 \theta} \log (1 + 2 \sin \theta) 4 \sin 2\theta  \,\mathrm{d}\theta \\ 
&= 4\int_0^{\pi/6} \theta \log (1 + 2 \sin \theta) \,\mathrm{d}\theta.
 \end{aligned}$$
This used the substitution $u^2 = t$, then $u = 2\sin \theta$ and 
then in the penultimate step 
$$ F'(z) = \frac{\arcsin \frac{\sqrt{z}}{2} }{\sqrt{1-z/4} \sqrt{z}} $$
to get $F'(4\sin^2 \theta) = \frac{\theta}{\sin 2\theta}$.
 Now factor
$$
1+2\sin\theta =
4\sin\!\left(\mfrac{\theta}{2}+\mfrac{\pi}{12}\right)
 \sin\!\left(\mfrac{\theta}{2}+\mfrac{5\pi}{12}\right) $$
to get
$$
\log(1+2\sin\theta) =
\log \Bigr(2\sin\bigl(\mfrac{\theta}{2}+\mfrac{\pi}{12}\bigr)\Bigr) +
\log \Bigr(2\sin\bigl(\mfrac{\theta}{2}+\mfrac{5\pi}{12}\bigr)\Bigr). $$

The standard Fourier series
$$
\log(2\sin \phi) =-\sum_{n=1}^\infty \frac{\cos(2n\phi )}{n}
$$ 
is uniformly convergent for $\frac{\theta}{2} + \frac{\pi}{12}$
and  $\frac{\theta}{2} + \frac{5\pi}{12}$ for $\theta \in [0, \frac{\pi}{6}]$.
Substituting for each summand in the right-hand side for $\log (1 + 2 \sin \theta)$ above we get
$$
\log(1+2\sin\theta)
=
-\sum_{n=1}^\infty
\frac{\cos\bigl(n(\theta+\frac{\pi}{6})\bigr)+\cos\bigl((\theta+\frac{5\pi}{6})n\bigr)}{n}.
$$
Using
$
\cos\bigl( n (\theta+\frac{\pi}{6}) \bigr) + \cos\bigl(n(\theta+\frac{5\pi}{6})\bigr)
= 2\cos\bigl(\frac{\pi}{3}n\bigr)\cos (\theta+\frac{\pi}{2})n
$
we obtain
$$
S
=
-8\sum_{n=1}^\infty \frac{\cos(\frac{\pi}{3}n)}{n}
\int_0^{\pi/6}\theta\cos \bigl( n (\theta+\frac{\pi}{2}) \bigr) \,\mathrm{d}\theta.
$$
A direct integration, which I verified in Mathematica, gives
$$
\int_0^{\pi/6}\theta\cos\bigl (n (\theta+\mfrac{\pi}{2}) \bigr) \,\mathrm{d}\theta
=
\frac{\pi}{6n}\sin\!\left(\frac{2n\pi}{3}\right)
+
\frac{\cos(\frac{2\pi}{3} n)-\cos( \frac{\pi}{2}n ) }{n^2}.
$$
Therefore using the uniform convergence mentioned above,
$$
S = -\frac{8\pi}{6}\sum_{n=1}^\infty \frac{\cos(\frac{\pi}{3}n)\sin(\frac{2\pi}{3} n)}{n^2} -8\sum_{n=1}^\infty \frac{a_n}{n^3},
$$
where
$$ a_n =  \cos \bigl( \mfrac{n\pi}{3} \bigr) \Bigl( 
\cos\bigl(\mfrac{2n\pi}{3}\bigr)-\cos\bigl(\mfrac{n\pi}{2}\bigr) \Bigr). $$
Since
$ \cos\bigl(\mfrac{n\pi}{3}\bigr)\sin \bigl(\mfrac{2n\pi}{3}\bigr)
= \frac{1}{2}\sin\bigl( \mfrac{n\pi}{3}\bigr) $
this becomes
$$ S =
-\frac{2\pi}{3}\sum_{n=1}^\infty \frac{\sin( \frac{\pi}{3} n)}{n^2}
-8\sum_{n=1}^\infty \frac{a_n}{n^3}.
$$

Now the sequence $(a_n)$ is periodic with period $12$.
The values $a_0,\ldots, a_{11}$ are 
$0,-\frac{1}{4},-\frac{1}{4},-1, \frac{3}{4}, -\frac{1}{4}, 2, -\frac{1}{4}, \frac{3}{4}, -1, -\frac{1}{4},
-\frac{1}{4}$. (Again I checked this with Mathematica.)
Hence
$$
a_n=
-\frac{1}{4}-\frac{3}{4} [3 \mid n] + 
[ 4 \mid n] + 3 [6 \mid n] - 3 [12 \mid n]
$$ 
in Iverson bracket notation.
It follows that 
$$
\sum_{n=1}^\infty \frac{a_n}{n^3} =
\Bigl( -\frac{1}{4} -\frac{3}{4\cdot 3^3} +\frac{1}{4^3} +\frac{3}{6^3} -\frac{3}{12^3}
\Bigr)\zeta(3)
= -\frac{1}{4}\zeta(3).
$$
We have shown that 
$$ S = 2\zeta(3)-\frac{2\pi}{3}\sum_{n=1}^\infty \frac{\sin(\frac{\pi}{3} n)}{n^2}. \qquad(\star) $$
It remains to evaluate 
$ C=\sum_{n=1}^\infty \frac{\sin(\frac{\pi}{3}n)}{n^2}$.
Using the values of $\sin \frac{\pi}{3}n$, which have period $6$, we get
$$
C = \frac{\sqrt{3}}{2}\sum_{m=0}^\infty
\bigl( \frac{1}{(6m+1)^2} +\frac{1}{(6m+2)^2} -\frac{1}{(6m+4)^2} -\frac{1}{(6m+5)^2} \bigr). $$
By one definition of the trigamma function $\psi_1$ we have
$\psi_1(z)=\sum_{m=0}^\infty \frac1{(m+z)^2}$,
so we may rewrite the previous equation for $C$ as
$$
C = \frac{\sqrt{3}}{72}
\Bigl(  \psi_1\bigl(\mfrac{1}{6} \bigr)  + \psi_1\bigl(\mfrac{1}{3} \bigr)  - \psi_1\bigl(\mfrac{2}{3} \bigr) 
-\psi_1\bigl(\mfrac{5}{6} \bigr)  \Bigr) $$
The <a href="https://www.johndcook.com/blog/gamma_identities/">trigamma duplication formula</a> states that
$$
\psi_1(z)+\psi_1\!\left(z+ \mfrac{1}{2} \right)=4\psi_1(2z). $$
Apply this with $z=\mfrac{1}{6}$ and $z = \mfrac{1}{3}$ to get
$ \psi_1\bigl( \mfrac{1}{6} \bigr) + \psi_1 \bigl( \mfrac{2}{3} \bigr) = 4 \psi_1 \bigl( \mfrac{1}{3} \bigr)$
and similarly
$ \psi_1\bigl( \mfrac{1}{3} \bigr) + \psi_1 \bigl( \mfrac{5}{6} \bigr) = 4 \psi_1 \bigl( \mfrac{2}{3} \bigr)$.
Hence
$$ \begin{aligned} C &=
\frac{\sqrt{3}}{72}
\Bigl( 4 \psi_1\bigl( \mfrac{1}{3} \bigr) - \psi_1 \bigl( \mfrac{2}{3} \bigr) 
+ \psi_1\bigl(\mfrac{1}{3} \bigr)  - \psi_1\bigl(\mfrac{2}{3} \bigr) 
-\psi_1\bigl(\mfrac{5}{6} \bigr) \Bigr) \\
&=
\frac{\sqrt{3}}{72}
\Bigl( 
5\psi_1\bigl(\mfrac{1}{3} \bigr) -2\psi_1\bigl(\mfrac{2}{3} \bigr) -\psi_1\bigl( \mfrac{5}{6} \bigr) \Bigr) 
\\ &=
\frac{\sqrt{3}}{72}
\Bigl( 
5\psi_1\bigl( \mfrac{1}{3} \bigr) - \mfrac{1}{2} \psi_1 \bigl( \mfrac{1}{3} \bigr) 
- \mfrac{1}{2} \psi_1 \bigl( \mfrac{5}{6} \bigr) - \psi_1 \bigl( \mfrac{5}{6} \bigr) 
\Bigr) \\
&=
\frac{\sqrt3}{48}
\bigl( 
3\psi_1\bigl( \mfrac{1}{3} \bigr) - \psi_1 \bigl( \mfrac{5}{6} \bigr) \bigr)
\end{aligned}
$$
Substituting this into the formula for $S$ marked ($\star$) above we arrive at
$$ S =
2\zeta(3) -\frac{2\pi}{3}C = 2\zeta(3)
-\frac{\pi\sqrt{3}}{24}\psi_1\bigl(\mfrac{1}{3} \bigr) 
+\frac{\pi\sqrt{3}}{72}\psi_1\bigl( \mfrac{5}{6} \bigr)
$$
as required.
